% Scheletro di lavoro, non testo definitivo.
% Le coordinate indicate sono piste di lettura da verificare sul testo greco prima della stesura.
% Metodo: testo originale, traduzione, ricostruzione argomentativa, analisi critica.
\chapter{La politica nelle opere tarde: il Politico e le Leggi}
\label{cap:platone:leggi}

% Problema: Che cosa distingue il vero politico da chi pretende di governare?
% Fonti da verificare: Politico 257a–268d.
\section{La definizione della competenza politica}

\subsection{La ricerca per divisione}

\subsection{La competenza del governo}

% Problema: Quale rapporto sussiste fra scienza politica, misura e legge?
% Fonti da verificare: Politico 268d–311c.
\section{La misura e il governo}

\subsection{Il mito e la correzione della definizione}

\subsection{La giusta misura, le costituzioni e le leggi}

% Problema: Come si può istituire un ordine politico governato dalle leggi?
% Fonti da verificare: Leggi I–III, VI; passi specifici da selezionare.
\section{La legislazione e la città delle Leggi}

\subsection{La finalità educativa delle leggi}

\subsection{L'assetto istituzionale della città}

% Problema: Le differenze fra i progetti politici indicano un'evoluzione o un mutamento del problema?
% Fonti da verificare: Repubblica V–VII; Politico; Leggi IV–XII.
\section{Continuità e differenze rispetto alla Repubblica}

\subsection{Filosofia e legislazione}

\subsection{Limiti delle letture evolutive}
